% GENERATED FILE. Edit canonical lessons, then run npm run generate:books.
% canonical-source-hash: fnv1a64:2fe7108467491e64
% canonical-lessons: ZH-C286-ganjue, ZH-W286-gan, ZH-C286-feeling, ZH-C286-duration, ZH-C286-duojiu, ZH-C286-yihan, ZH-C286-practice

\chapter{How It Felt}
\label{ch:zh-how-it-felt}
\hlchaptermodality{286}

\begin{chapteropening}
\textbf{By the end of this chapter:} \emph{I can say how something felt and how long it lasted.}

Saying how something felt with \zh{我感觉} and a feeling, how long it lasted with the length after the verb, how long and what a pity, built on one new written character: \zh{感}.
\end{chapteropening}

\section[\texorpdfstring{\zh{感觉} (gǎnjué)}{感觉 (gǎnjué)}]{\zh{感觉} (gǎnjué) — to feel; a feeling}

\label{lesson:ZH-C286-ganjue}



\begin{quote}
A new chapter, and it opens on how things felt.

\begin{itemize}
\raggedright
  \item \emph{Recall:} say \zh{我的车票不见了} — \emph{wǒ de chēpiào bú jiàn le} — \textbf{R1}, the chapter you have finished
  \item \emph{Recall:} say \emph{gǎndào}, to feel — \textbf{R3}, durable
  \item \emph{Recall:} say \zh{开心} — \emph{kāixīn}, happy — \textbf{R3}
\end{itemize}
\end{quote}

\subsection*{You'll want to know: \zh{感觉}}
\textbf{\zh{感觉}} — \emph{gǎnjué} — \textquotedblleft{}to feel\textquotedblright{}, how something seems to you; \textquotedblleft{}a feeling\textquotedblright{}.

\begin{quote}
\textbf{\zh{我感觉很好}} \emph{wǒ gǎnjué hěn hǎo} — I feel fine.
\end{quote}

\begin{quote}
\emph{wǒ gǎnjué tā bù gāoxìng} — I have a feeling he's not happy.
\end{quote}

The next lesson puts the first character in your hand. The second is the \zh{觉} of \emph{juéde}.

\begin{cousinweb}
\textbf{\zh{感觉}} (\emph{gǎnjué}) is built from \zh{感} and \zh{觉}, both inherited from Middle Chinese. \zh{感} is \emph{komX}, glossed 'to move, to stir'. The Guangyun lists \zh{觉} under two readings: \emph{kaewH}, glossed as sleep, and \emph{kaewk}, glossed 'to understand' and 'to wake'; jué is the regular modern sound of \emph{kaewk} (Guangyun, in Baxter's transcription). 'To feel' for the pair is the word's own sense.
\end{cousinweb}

\begin{grammarlens}[title={What you've built}]
\emph{gǎndào} and \zh{感觉} both report a feeling; \zh{感觉} is also how something seems to you: \zh{我感觉很好}.

\emph{juéde}, as in \zh{我觉得}, is what you think about it.
\end{grammarlens}

\subsection*{Your turn}
\begin{itemize}
\raggedright
  \item \emph{Say it:} \emph{gǎnjué}
  \item \emph{Say it:} \emph{gǎnjué}, once more
  \item \emph{Say it:} \zh{我感觉很好} — \emph{wǒ gǎnjué hěn hǎo} — I feel fine.
  \item \emph{Say it:} \emph{wǒ gǎnjué tā bù gāoxìng} — I have a feeling he's not happy.
  \item \emph{Recall:} say \zh{我觉得}, then \zh{我感觉}: what I think, then how it feels
\end{itemize}

\begin{tcolorbox}[breakable,colback=teal!4,colframe=teal!35!black,title={Before you move on}]
What does \zh{感觉} mean? (\textquotedblleft{}To feel; a feeling\textquotedblright{}.) Which word gives an opinion instead? (\textbf{\zh{觉得}}.)
\end{tcolorbox}
\section[\texorpdfstring{\zh{感} (gǎn)}{感 (gǎn)}]{\zh{感} — a frame with \zh{口} inside, then \zh{心}}

\label{lesson:ZH-W286-gan}



\begin{quote}
The word first, then the character inside it.

\begin{itemize}
\raggedright
  \item \emph{Recall:} say \emph{gǎnjué} — \textbf{R1}, one lesson back, before it decays
\end{itemize}

\begin{quote}
\zh{我} \zh{口} \zh{心}
\end{quote}

\emph{Read it:} each one aloud, with nothing said first — \emph{wǒ}, \emph{kǒu}, \emph{xīn}
\end{quote}

\subsection*{Script — a frame with \zh{口} inside, then \zh{心}}
\hlblockfigure{\detokenize{figures/ZH-W286-gan-filmstrip.pdf}}{How \zh{感} is written, stroke by stroke}

\begin{quote}
\zh{感}
\end{quote}

\textbf{Two pieces: a frame with \zh{口} inside it, then \zh{心}.}

\textbf{The frame, nine strokes}: a long horizontal at the top; a left-falling stroke down the left side; a short horizontal under the top; \zh{口}, three strokes, small, under that; then the long slanting hook down to the lower right, as in \zh{我}; a short left-falling stroke across it; and last the dot at the top right.

\textbf{\zh{心}, four strokes}, at the foot: the left dot, the curve along the foot that hooks up, the middle dot, the right dot.

9 + 4 = thirteen.

\subsection*{Your turn}
\begin{itemize}
\raggedright
  \item \emph{Write it:} \zh{感}, all thirteen, with the model in front of you
  \item \emph{Cover:} write it again from memory, and count the pieces, not the strokes
\end{itemize}

\begin{tcolorbox}[breakable,colback=teal!4,colframe=teal!35!black,title={Before you move on}]
How many strokes in \zh{感}? (\textbf{Thirteen}.) Where have you written the slanting hook? (\textbf{In \zh{我}}.) What is at the foot? (\textbf{\zh{心}}.)

Source: \href{https://raw.githubusercontent.com/chanind/hanzi-writer-data/68d10a4b21150cae5e1ebbd223eed289cf32d90c/data/\%E6\%84\%9F.json}{Hanzi Writer Data \zh{感}, pinned stroke paths and medians}.
\end{tcolorbox}
\section[\texorpdfstring{\zh{我感觉}… (wǒ gǎnjué …)}{我感觉… (wǒ gǎnjué …)}]{\zh{我感觉}… — how it felt}

\label{lesson:ZH-C286-feeling}



\begin{quote}
A word you have, and the feelings that go after it.

\begin{itemize}
\raggedright
  \item \emph{Read it:} \zh{感觉}, aloud — \emph{gǎnjué}
  \item \emph{Recall:} say \emph{nánguò}, sad — \textbf{R3}, durable
  \item \emph{Recall:} say \emph{jǐnzhāng}, nervous — \textbf{R3}
  \item \emph{Recall:} say \emph{shīwàng}, disappointed — \textbf{R3}
\end{itemize}
\end{quote}

\begin{grammarlens}[title={Grammar Lens: \zh{我感觉} and a feeling}]
\begin{quote}
\textbf{\zh{我感觉} and a feeling word says how it felt; the feeling words are ones you know.}
\end{quote}

\begin{quote}
\textbf{\zh{那天我感觉很开心}} \emph{nà tiān wǒ gǎnjué hěn kāixīn} — I felt really happy that day. \emph{shàngkè de shíhou wǒ gǎnjué hěn jǐnzhāng} — I felt nervous in class.
\end{quote}

The feeling takes \zh{很}, as it does after \zh{我}: \zh{我感觉很开心}. Any feeling you know fits: \emph{gāoxìng}, \emph{nánguò}, \emph{jǐnzhāng}, \emph{jīdòng}, \emph{shīwàng}.

\zh{感觉} is for how it felt; \zh{觉得} is for what you think. \zh{我觉得他说得对} is an opinion; \zh{我感觉很好} is a feeling.
\end{grammarlens}

\subsection*{Your turn}
Say these aloud:

\begin{itemize}
\raggedright
  \item \zh{那天我感觉很开心} — \emph{nà tiān wǒ gǎnjué hěn kāixīn} — I felt really happy that day.
  \item \emph{shàngkè de shíhou wǒ gǎnjué hěn jǐnzhāng} — I felt nervous in class.
  \item how you felt on your last birthday, with \zh{我感觉}
\end{itemize}

\begin{tcolorbox}[breakable,colback=teal!4,colframe=teal!35!black,title={Before you move on}]
What comes after \zh{我感觉}? (\textbf{A feeling, with \zh{很}}.) Which word gives an opinion instead? (\textbf{\zh{觉得}}.)
\end{tcolorbox}
\section[\texorpdfstring{\zh{学了两年} (xué le liǎng nián)}{学了两年 (xué le liǎng nián)}]{\zh{学了两年} — how long it lasted}

\label{lesson:ZH-C286-duration}



\begin{quote}
Numbers you know, and a new place to put them.

\begin{itemize}
\raggedright
  \item \emph{Read it:} \zh{那天我感觉很开心}, aloud — \emph{nà tiān wǒ gǎnjué hěn kāixīn}
  \item \emph{Recall:} say \zh{年} — \emph{nián}, a year — \textbf{R3}, durable
  \item \emph{Recall:} say \emph{hǎojiǔ}, a long time — \textbf{R3}
  \item \emph{Recall:} say \zh{多大} — \emph{duō dà}, how old — \textbf{R3}
\end{itemize}
\end{quote}

\begin{grammarlens}[title={Grammar Lens: the length after the verb}]
\begin{quote}
\textbf{How long something lasted goes after the verb and its \zh{了}; with an object, the object comes after the length.}
\end{quote}

\begin{quote}
\textbf{\zh{我学了两年汉语}} \emph{wǒ xué le liǎng nián Hànyǔ} — I studied Chinese for two years. \emph{tā zài Zhōngguó zhù le yì nián} — He lived in China for a year.
\end{quote}

When it happened goes before the verb: \zh{去年}. How long it lasted goes after: \zh{两年}.

To ask, put \zh{多} before \zh{长}, as you put it before \zh{大} in \zh{多大}: \textbf{\zh{你学了多长时间}} \emph{nǐ xué le duō cháng shíjiān}?
\end{grammarlens}

\subsection*{Your turn}
Say these aloud:

\begin{itemize}
\raggedright
  \item \zh{我学了两年汉语} — \emph{wǒ xué le liǎng nián Hànyǔ} — I studied Chinese for two years.
  \item \emph{tā zài Zhōngguó zhù le yì nián} — He lived in China for a year.
  \item \zh{你学了多长时间} — \emph{nǐ xué le duō cháng shíjiān} — How long did you study?
  \item how long you have been learning Chinese, with \zh{了} and the length
\end{itemize}

\begin{tcolorbox}[breakable,colback=teal!4,colframe=teal!35!black,title={Before you move on}]
Where does the length go? (\textbf{After the verb and \zh{了}}.) And the object? (\textbf{After the length: \zh{两年汉语}}.)
\end{tcolorbox}
\section[duōjiǔ]{duōjiǔ — how long}

\label{lesson:ZH-C286-duojiu}



\begin{quote}
Before the new one: last lesson's question, and an answer you already have.

\begin{itemize}
\raggedright
  \item \emph{Read it:} \zh{你学了多长时间}, aloud — \emph{nǐ xué le duō cháng shíjiān}
  \item \emph{Recall:} say \emph{hǎojiǔ}, a long time — \textbf{R3}, durable
\end{itemize}
\end{quote}

\subsection*{You'll want to know: duōjiǔ}
\emph{duōjiǔ} — \textquotedblleft{}how long\textquotedblright{}, of time.

\begin{quote}
\emph{nǐ xué le duōjiǔ} — How long did you study?
\end{quote}

\begin{quote}
\emph{nǐ zài Zhōngguó zhù le duōjiǔ} — How long did you live in China?
\end{quote}

Its second character comes later; for now it is a word for the ear.

\begin{cousinweb}
\emph{duōjiǔ} is built from \zh{多} 'much' and jiǔ 'long-lasting', both inherited from Middle Chinese, \emph{ta} and \emph{kjuwX} (Guangyun, in Baxter's transcription), which glosses the first 'many' and the second 'lasting long'.
\end{cousinweb}

\begin{grammarlens}[title={What you've built}]
\emph{duōjiǔ} is the short form of \zh{多长时间}, and it takes the same place, after the verb.

\emph{hǎojiǔ} answers it: a long time.
\end{grammarlens}

\subsection*{Your turn}
\begin{itemize}
\raggedright
  \item \emph{Say it:} \emph{duōjiǔ}
  \item \emph{Say it:} \emph{duōjiǔ}, once more
  \item \emph{Say it:} \emph{nǐ xué le duōjiǔ} — How long did you study?
  \item \emph{Say it:} \emph{nǐ zài Zhōngguó zhù le duōjiǔ} — How long did you live in China?
  \item \emph{Recall:} say \zh{多长时间}, then \emph{duōjiǔ}: the long question, then the short one
\end{itemize}

\begin{tcolorbox}[breakable,colback=teal!4,colframe=teal!35!black,title={Before you move on}]
What does \emph{duōjiǔ} mean? (\textquotedblleft{}How long\textquotedblright{}.) Which longer question means the same? (\textbf{\zh{多长时间}}.)
\end{tcolorbox}
\section[yíhàn]{yíhàn — a pity}

\label{lesson:ZH-C286-yihan}



\begin{quote}
Before the new one: last lesson's question, and two words for when things go wrong.

\begin{itemize}
\raggedright
  \item \emph{Recall:} say \emph{duōjiǔ}, how long — \textbf{R1}, one lesson back
  \item \emph{Recall:} say \emph{hòuhuǐ}, to regret — \textbf{R3}, durable
  \item \emph{Recall:} say \emph{shīwàng}, disappointed — \textbf{R3}
\end{itemize}
\end{quote}

\subsection*{You'll want to know: yíhàn}
\emph{yíhàn} — \textquotedblleft{}a pity; regrettable\textquotedblright{}.

\begin{quote}
\emph{tài yíhàn le} — What a pity!
\end{quote}

\begin{quote}
\emph{wǒ méi qù, hěn yíhàn} — I didn't go, which is a shame.
\end{quote}

Its characters come later; for now it is a word for the ear.

\begin{cousinweb}
\emph{yíhàn} is built from yí 'to leave behind' and hàn 'resentment', both inherited from Middle Chinese. The Guangyun lists the first under two readings, \emph{ywij}, glossed 'to lose' and 'to leave', and \emph{ywijH}, glossed 'to give'; the second is \emph{homH}, glossed 'resentment' (Guangyun, in Baxter's transcription). 'A pity' for the pair is the modern word's own sense.
\end{cousinweb}

\begin{grammarlens}[title={What you've built}]
\emph{hòuhuǐ} is regretting what you did yourself; \emph{yíhàn} is a pity, whoever was to blame.

It takes \zh{很} or \zh{太} … \zh{了}, as a feeling does: \emph{tài yíhàn le}.
\end{grammarlens}

\subsection*{Your turn}
\begin{itemize}
\raggedright
  \item \emph{Say it:} \emph{yíhàn}
  \item \emph{Say it:} \emph{yíhàn}, once more
  \item \emph{Say it:} \emph{tài yíhàn le} — What a pity!
  \item \emph{Say it:} \emph{wǒ méi qù, hěn yíhàn} — I didn't go, which is a shame.
  \item \emph{Recall:} say \emph{hòuhuǐ}, then \emph{yíhàn}: my fault, then a pity
\end{itemize}

\begin{tcolorbox}[breakable,colback=teal!4,colframe=teal!35!black,title={Before you move on}]
What does \emph{yíhàn} mean? (\textquotedblleft{}A pity\textquotedblright{}.) How do you say what a pity? (\emph{tài yíhàn le}.)
\end{tcolorbox}
\section[Practice]{\zh{我感觉很好}}

\label{lesson:ZH-C286-practice}



\begin{quote}
The chapter closes here, with someone back from a year away.

\begin{itemize}
\raggedright
  \item \emph{Recall:} say \emph{yíhàn}, a pity — \textbf{R1}, one lesson back
  \item \emph{Recall:} say \emph{duōjiǔ}, how long — \textbf{R2}
  \item \emph{Recall:} say \zh{我学了两年汉语} — \emph{wǒ xué le liǎng nián Hànyǔ} — \textbf{R2}
\end{itemize}
\end{quote}

\subsection*{The exchange}
\begin{quote}
\textbf{A:} \zh{你去过中国吗}
\end{quote}

\begin{quote}
\emph{nǐ qù guo Zhōngguó ma} — Have you been to China?
\end{quote}

\begin{quote}
\textbf{B:} \zh{去过} \zh{我在中国学了一年汉语}
\end{quote}

\begin{quote}
\emph{qù guo, wǒ zài Zhōngguó xué le yì nián Hànyǔ} — I have. I studied Chinese there for a year.
\end{quote}

\begin{quote}
\textbf{A:} \emph{nǐ gǎnjué zěnmeyàng} — How did it feel?
\end{quote}

\begin{quote}
\textbf{B:} \zh{我感觉很好}
\end{quote}

\begin{quote}
\emph{wǒ gǎnjué hěn hǎo} — It felt good.
\end{quote}

\begin{quote}
\textbf{A:} \emph{nǐ zài Zhōngguó zhù le duōjiǔ} — How long did you live there?
\end{quote}

\begin{quote}
\textbf{B:} \zh{一年}
\end{quote}

\begin{quote}
\emph{yì nián} — A year.
\end{quote}

\begin{quote}
\textbf{A:} \zh{你想再去吗}
\end{quote}

\begin{quote}
\emph{nǐ xiǎng zài qù ma} — Do you want to go again?
\end{quote}

\begin{quote}
\textbf{B:} \emph{xiǎng, méiyǒu shíjiān, tài yíhàn le} — I do, but there's no time. What a pity.
\end{quote}

A asks with \zh{过}, then with \zh{感觉}, then with \emph{duōjiǔ}; B answers each in as few words as will do: \zh{去过}, \zh{一年}.

\begin{tcolorbox}[breakable,skin=enhanced,colback=orange!6,colframe=orange!45!black,boxrule=0.5pt,arc=1mm,left=6pt,right=6pt,top=4pt,bottom=4pt,fonttitle=\bfseries,title={Read this}]
\emph{Read it:} the exchange through once without stopping, then answer the question below
\end{tcolorbox}

\begin{grammarlens}[title={What you've built — how it was}]
How it felt, how long it lasted, and how you feel about it now:

\begin{itemize}
\raggedright
  \item \textbf{\zh{那天我感觉很开心}} \emph{nà tiān wǒ gǎnjué hěn kāixīn} — how it felt.
  \item \textbf{\zh{我学了两年汉语}} \emph{wǒ xué le liǎng nián Hànyǔ} — how long.
  \item \emph{nǐ xué le duōjiǔ?} — asking how long.
  \item \emph{tài yíhàn le} — a pity.
\end{itemize}

\textbf{Listen.} \emph{Nǐ zài nàr zhù le duōjiǔ?}

\textbf{Speak.} Tell someone about a place you stayed: how long, and how it felt.

\emph{Write it:} \zh{感觉} — the frame with \zh{口}, then \zh{心}
\end{grammarlens}

\subsection*{Your turn}
\begin{itemize}
\raggedright
  \item \emph{Say it:} \zh{我在中国学了一年汉语} \zh{感觉很好} — \emph{wǒ zài Zhōngguó xué le yì nián Hànyǔ, gǎnjué hěn hǎo} — I studied Chinese in China for a year, and it felt good.
  \item \emph{Say it:} \zh{你学了多长时间} — \emph{nǐ xué le duō cháng shíjiān} — How long did you study?
  \item \emph{Say it:} the whole exchange, both parts
  \item \emph{Write it:} \zh{我感觉很好}
\end{itemize}

\begin{tcolorbox}[breakable,colback=teal!4,colframe=teal!35!black,title={Before you move on}]
Say you feel fine. (\textbf{\zh{我感觉很好}}.) Say you studied Chinese for two years. (\textbf{\zh{我学了两年汉语}}.) Say what a pity. (\emph{tài yíhàn le}.)
\end{tcolorbox}
